\section{Polynomial differential operators and the Toda radial part}
\label{rt-section}

We now construct two of the actual differential operators used
in the Fourier argument.  Their coefficients are polynomial on
the full unitary group.  We also prove their precise restrictions
to the radial cell, including the scalar conjugation and the
conjugate-coordinate operator.  The proof of the radial identity
uses a last-column expansion of a Capelli determinant and the
two unipotent equivariances; no general theorem about radial
parts is required.  These are the operators denoted by
the weighted isotropic Capelli minors in the source
\cite{alpo-uniform-source}.

\subsection{Right derivatives and the two polynomial operators}

Keep the Hermitian space and notation of Section~\ref{rl-section}.
In particular $V=U\oplus W\oplus U^\vee$,
$\dim U=r$, and $\dim W=k+1$.
Fix any ordering of its chosen Witt basis.
For $X$ in the real Lie algebra of $\mathbf G=U(V,h)$ put
\[
 (R_XF)(g)=\left.\frac{d}{dt}\right|_{t=0}F(g\exp(tX)).
\]
For an arbitrary complex matrix $A$ define the complexified
right derivative by
\begin{equation}\label{rt-complex-derivative}
 R_AF=R_{(A-HA^*H)/2}F
            +iR_{(A+HA^*H)/(2i)}F.
\end{equation}
Both matrices on the right belong to the real Lie algebra.
This is complex linear in $A$, and
$[R_A,R_B]=R_{[A,B]}$.  Indeed the real right derivatives
are the left-invariant vector fields with tangent vectors
$gA$, and differentiating these vector fields gives this
bracket; complexification preserves it.

Write $E_{ab}=R_{e_ae_b^*}$, where $e_b^*$ here means
the coordinate covector, not the Hermitian dual.
For ordered lists $I=(i_1,\ldots,i_m)$ and
$J=(j_1,\ldots,j_m)$ of distinct basis indices define
\begin{equation}\label{rt-minor}
 Q_{I,J}(u)=
 \sum_{\sigma\in S_m}\operatorname{sgn}(\sigma)
 \prod_{b=1}^{m}
 \bigl((u-b+1)\delta_{i_{\sigma(b)},j_b}
                         +E_{i_{\sigma(b)},j_b}\bigr).
\end{equation}
Products act on a function with the rightmost factor first.
The empty minor is the identity.  These are the minors
of Theorem~\ref{rp-capelli}.

Set $I=(e_1,\ldots,e_r)$ and
$I^\vee=(e_1^\vee,\ldots,e_r^\vee)$.  If
$J=(j_1,\ldots,j_r)$ runs over the increasing $r$-element
subsets of the chosen basis, put
\begin{align}
 \beta_J(g)&=\det\bigl(h(ge_{j_b},e_a)\bigr)_{a,b=1}^{r},
                                                        \label{rt-beta}\\
 \alpha_J(g)&=\det\bigl((g^{-1}e_a)_{j_b}\bigr)_{a,b=1}^{r}.
                                                        \label{rt-alpha}
\end{align}
The transpose in the second displayed determinant does
not change its value.  For $Y\in\mathbb C$ set
\begin{align}
 \mathcal A_YF(g)
   &=\sum_J\beta_J(g)\,
       Q_{I,J}\left(Y+\frac{2r+k}{2}\right)F(g),
                                                   \label{rt-A}\\
 \overline{\mathcal A}_YF(g)
   &=\sum_J\alpha_J(g)\,
       Q_{J,I^\vee}\left(Y-\frac{2r+k}{2}+r-1\right)F(g).
                                                   \label{rt-Av}
\end{align}
The bar in $\overline{\mathcal A}$ is a label for the
conjugate-coordinate operator; it does not conjugate the
argument $F$ or the parameter $Y$.
Changing the order of one summation list changes both
its scalar determinant and its quantum minor by the
same sign.  Hence any fixed rule for ordering the subsets
gives the same operators.  This explains the different
ordering of the dual indices used in the source.
For $r=0$ both operators are the identity.

\begin{proposition}\label{rt-polynomial}
The operators \eqref{rt-A} and \eqref{rt-Av} have
differential order at most $r$ and polynomial coefficients
of degree $r$ in $g$ and $\overline g$, respectively.
They are smooth on the whole group, including the zeros
of every intermediate minor.

Every derivative of a scalar coefficient, of any fixed
length $m$ in left or right Lie derivatives with directions
of norm at most one, is bounded by
$C_m(1+\|g\|)^r$.  The constant is independent of $g$
and of those directions.  Their formal adjoints for
Haar measure also have polynomial coefficients of degree
at most $r$ and order at most $r$.

Both the operators and their formal adjoints preserve the
class $\mathscr S_\varepsilon$ of
Corollary~\ref{rl-class-preserved}.
\end{proposition}

\begin{proof}
The first coefficients in \eqref{rt-beta} are polynomial
of degree $r$.  On $\mathbf G$ one has $g^{-1}=Hg^*H$,
so the second coefficients are polynomial of degree $r$
in $\overline g$.  Each term in \eqref{rt-minor} is a
constant linear combination of words of length at most
$r$ in right derivatives.  Thus no division occurs in
either operator.

A left or right Lie derivative sends a matrix entry of
$g$ to a linear combination of its entries, and likewise
for conjugate entries.  On the finite-dimensional space
of polynomials of degree at most $r$, this derivative
therefore acts by a linear map whose coefficients are
bounded when the direction has norm at most one.
Iterating $m$ times bounds the sum of the absolute
coefficients by a fixed constant $C_m$.
Evaluating the resulting polynomials proves the
stated bound.

For clarity, Haar measure on this group is unimodular.
The real bilinear form
$B(X,Z)=\operatorname{Re}\operatorname{tr}(XZ)$ on its
Lie algebra is nondegenerate and invariant under the
adjoint action of every group element.  Nondegeneracy
follows by complexifying: the complexified Lie algebra
is the full matrix algebra, whose trace pairing is
nondegenerate.  An adjoint transformation preserves
$B$, so its real determinant has absolute value one.
The absolute volume form defined at the identity
therefore extends to a left and right invariant
measure.  Integration by parts for compactly supported
functions gives $R_X^*=-R_X$ for real $X$.
Taking adjoints of \eqref{rt-A} and \eqref{rt-Av}
reverses each word and conjugates its coefficients.
Leibniz's rule differentiates those polynomials but
does not increase their degree.

Finally every derivative of an output is a finite sum
of a polynomially bounded coefficient times a derivative
of the input.  Arbitrarily fast polynomial decay of
all input derivatives implies the same for the output.
A differential operator does not enlarge support.
This proves preservation of $\mathscr S_\varepsilon$
for the operators and their adjoints.
\end{proof}

\subsection{Equivariance on the full group}

Let $N,K_z,K_{z'}$ be as in
Lemma~\ref{rl-cofactor-covariance}.
The following identities hold for arbitrary smooth
functions, before imposing a character on $N$.

\begin{proposition}\label{rt-full-covariance}
If $s_L\in K_zN$ and $s_R\in K_{z'}N$, then
\begin{align*}
 \mathcal A_Y\bigl(F(s_L\,\cdot\,s_R)\bigr)(g)
       &=(\mathcal A_YF)(s_Lgs_R),\\
 \overline{\mathcal A}_Y\bigl(F(s_L\,\cdot\,s_R)\bigr)(g)
       &=(\overline{\mathcal A}_YF)(s_Lgs_R).
\end{align*}
The same assertions hold for vector-valued functions
component by component.
\end{proposition}

\begin{proof}
Both $s_L$ and $s_R$ preserve $U$ and act on it with
determinant one.  They also preserve $U^\perp$, and
the induced action on $V/U^\perp$, in the basis
given by $I^\vee$, has determinant one.

Left translation commutes with right derivatives.
For the first operator,
\[
 h(s_Lge_j,e_a)=h(ge_j,s_L^{-1}e_a).
\]
The basis change on $U$ has determinant one, so every
$r$-row determinant in \eqref{rt-beta} is unchanged.
For the second operator,
$\bigwedge^r(s_L^{-1})$ fixes
$e_1\wedge\cdots\wedge e_r$, so
$\alpha_J(s_Lg)=\alpha_J(g)$.
This proves left covariance.

For right translation by $s$, complexified derivatives
transform by
\[
 R_A\bigl(F(\,\cdot\,s)\bigr)(g)
          =(R_{s^{-1}As}F)(gs).
\]
The exterior covariance of the quantum minors gives
\begin{equation}\label{rt-right-minor}
 Q'_{I,J}(u)
    =\sum_{K,L}\det(s^{-1})_{K,I}\,
                   \det(s)_{J,L}\,Q_{K,L}(u),
\end{equation}
where $Q'$ is formed using $A\mapsto R_{s^{-1}As}$.
This is precisely the all-minor covariance in
Theorem~\ref{rp-capelli}; no coefficient is
differentiated in this constant change of basis.

Since $\bigwedge^r(s^{-1})$ fixes $\bigwedge^rU$,
the first determinant on the right of
\eqref{rt-right-minor}, when $I$ is the fixed list
of $U$-indices, equals one for $K=I$ and zero otherwise.
Cauchy--Binet then gives
\[
 \sum_J \beta_J(g)\det(s)_{J,L}=\beta_L(gs).
\]
This proves right covariance of $\mathcal A_Y$.

For the second operator, Cauchy--Binet gives instead
\[
 \sum_J\det(s^{-1})_{K,J}\alpha_J(g)=\alpha_K(gs).
\]
The remaining coefficients in
\eqref{rt-right-minor} are $\det(s)_{I^\vee,L}$.
The last $r$ coordinate rows of $s$ vanish on
$U^\perp$ and have determinant one on the quotient.
Consequently these coefficients are one for
$L=I^\vee$ and zero otherwise.  This proves right
covariance of $\overline{\mathcal A}_Y$.
The argument is unaffected by any fixed reordering
of the dual basis.
\end{proof}

\subsection{The two unipotent characters and elementary derivative rules}

For $n\in N$, define
\[
 \ell_z(n)=
 \sum_{i=1}^{r-1}(ne_{i+1})_{e_i}+(nz)_{e_r},
 \qquad
 \psi_z(n)=\exp\bigl(\ell_z(n)-\overline{\ell_z(n)}\bigr).
\]
For $r=0$ set $\ell_z=0$.
The same definition with $z'$ gives $\psi_{z'}$.
These are the actual normalizations of the source's
characters.  Only their values on the roots inside
$U$ and $U^\vee$ enter the proof below.

Let $D$ be an open subset of $\mathbf G$ invariant
under multiplication by $N$ on both sides.
Let $F:D\to E$ be smooth, where $E$ is a
finite-dimensional complex vector space, and suppose
\begin{equation}\label{rt-whittaker}
 F(n_Lgn_R)=\psi_z(n_L)\psi_{z'}(n_R)F(g).
\end{equation}
Fix $m\in U(W)$ and consider the points $b(c)m\in D$,
with all $c_i\ne0$, as in Lemma~\ref{rl-cell}.
Write
\[
 F_m(c)=F(b(c)m),\qquad
 p_i=c_i\partial_{c_i},\qquad
 \bar p_i=\overline c_i\partial_{\overline c_i}.
\]
Wirtinger derivatives are meant in these formulas.

\begin{lemma}\label{rt-root-rules}
For the ordered indices $e_1,\ldots,e_r$,
right differentiation on a function satisfying
\eqref{rt-whittaker} gives
\begin{align}
 E_{e_i,e_j}F&=\delta_{j,i+1}F
                                      &&(i<j),\label{rt-right-rule}\\
 (E_{e_i,e_j}F)(b(c)m)
       &=\delta_{i,j+1}\frac{c_i}{c_j}F_m(c)
                                      &&(i>j),\label{rt-left-rule}\\
 (E_{e_i,e_i}F)(b(c)m)&=p_iF_m(c).
                                                   \label{rt-torus-rule}
\end{align}
Equation \eqref{rt-left-rule} holds also with $F$
replaced by any word of right derivatives applied
to $F$.

For the reversed dual list
$v_i=e_{r+1-i}^\vee$, the corresponding rules are
\begin{align}
 E_{v_i,v_j}F&=\delta_{j,i+1}F&&(i<j),\label{rt-dual-right}\\
 (E_{v_i,v_j}F)(b(c)m)
   &=\delta_{i,j+1}
       \frac{\overline c_{r+1-j}}{\overline c_{r+1-i}}
         F_m(c)&&(i>j),\label{rt-dual-left}\\
 (E_{v_i,v_i}F)(b(c)m)
   &=-\bar p_{r+1-i}F_m(c).\label{rt-dual-torus}
\end{align}
The dual lower-root rule again holds after any
word of right derivatives.
\end{lemma}

\begin{proof}
For $i<j$, the real root subgroup inside $N$ is
\[
 n_{ij}(w)=
 \exp\bigl(w e_ie_j^*-\overline w\,
                         e_j^\vee(e_i^\vee)^*\bigr).
\]
Its character is $\exp(w-\overline w)$ when
$j=i+1$ and is one otherwise.
Complexifying its two real tangent directions as
in \eqref{rt-complex-derivative} amounts to taking
$\partial_w$ at zero.  The right character law
therefore gives \eqref{rt-right-rule}.

At $g=b(c)m$ one has, for $i>j$,
\[
 \operatorname{Ad}(g)(e_ie_j^*)
          =\frac{c_i}{c_j}
                      e_i^\vee(e_j^\vee)^*.
\]
The matrix on the right is an upper root in the
dual part of $N$.  Its character derivative is
one when $i=j+1$ and zero otherwise.  For example,
if its root parameter is $w$, the corresponding
upper entry in the $U$-block is $-\overline w$,
whose contribution to
$\ell-\overline\ell$ is again $w-\overline w$.
Transferring the right derivative at $g$ to this
left derivative proves \eqref{rt-left-rule}.
Every right derivative commutes with left
translations.  Thus every right-differentiated
function retains the left character law, which
proves the asserted strengthened version.

The two real directions in
\eqref{rt-complex-derivative} for $e_ie_i^*$
scale $c_i$ respectively by $e^{t/2}$ and
$e^{-it/2}$.  Their complex combination is
$c_i\partial_{c_i}$, proving \eqref{rt-torus-rule}.

In the reversed dual ordering, the upper root
$e_a^\vee(e_b^\vee)^*$ has $a>b$ and has
the same character derivative, equal to one
precisely when $a=b+1$.  For a lower root in
that ordering, conjugation by $b(c)m$ gives
\[
 \operatorname{Ad}(b(c)m)
       \bigl(e_a^\vee(e_b^\vee)^*\bigr)
           =\frac{\overline c_b}{\overline c_a}
                          e_ae_b^* .
\]
This proves \eqref{rt-dual-right} and
\eqref{rt-dual-left}.  The real torus directions
for $e_i^\vee(e_i^\vee)^*$ scale $c_i$ by
$e^{-t/2}$ and $e^{-it/2}$; their complex
combination is $-\bar c_i\partial_{\bar c_i}$.
This proves the last identity.
\end{proof}

\subsection{An elementary determinant recurrence}

The following argument is the essential radial calculation.
It explains why only adjacent roots survive.

\begin{lemma}\label{rt-cofactor-reduction}
Let $E_{ab}$ satisfy the $\mathfrak{gl}_s$ commutator
relations, and let $s\geq2$.
In its universal enveloping algebra, put
\[
 C_s(u)=Q_{(1,\ldots,s-2,s),(1,\ldots,s-1)}(u).
\]
Then
\begin{equation}\label{rt-cofactor-ideal}
 C_s(u)=E_{s,s-1}Q_{(1,\ldots,s-2),(1,\ldots,s-2)}(u)
                  +\sum_{j=1}^{s-2}E_{sj}V_j
\end{equation}
for some elements $V_j$ of the enveloping algebra.
No assertion about the values of the $V_j$ is needed.
\end{lemma}

\begin{proof}
In each term of the determinant, there is exactly
one factor from row $s$, necessarily an off-diagonal
factor $E_{sj}$.  If $j<s-1$, move this factor to
the far left.  For a factor from an earlier column
$b<j$ one has
\[
 \bigl((u-b+1)\delta_{ab}+E_{ab}\bigr)E_{sj}
  =E_{sj}\bigl((u-b+1)\delta_{ab}+E_{ab}\bigr)
                                  -\delta_{aj}E_{sb}.
\]
Every new last-row generator has an even smaller
column index.  Repeating this identity expresses
the term as a sum of elements with leftmost factor
$E_{s\ell}$, where $\ell\leq j<s-1$.

If $j=s-1$, every preceding row index is at most
$s-2$.  The same commutator is then zero.
Moving $E_{s,s-1}$ to the far left leaves exactly
the determinant on the first $s-2$ rows and
columns, with parameter $u$ and positive sign.
Summing all terms proves \eqref{rt-cofactor-ideal}.
\end{proof}

\begin{proposition}\label{rt-raw-toda}
For $0\leq s\leq r$, let $K_s(u)$ be the
coordinate differential operator defined by
\[
 K_0(u)=1,\qquad K_1(u)=u+p_1,
\]
and, for $s\geq2$,
\begin{equation}\label{rt-K-rec}
 K_s(u)=(u-s+1+p_s)K_{s-1}(u)
                      -\frac{c_s}{c_{s-1}}K_{s-2}(u).
\end{equation}
For every $F$ satisfying \eqref{rt-whittaker},
\begin{equation}\label{rt-raw-identity}
 \bigl(Q_{(e_1,\ldots,e_s),(e_1,\ldots,e_s)}(u)F\bigr)
                (b(c)m)=K_s(u)F_m(c).
\end{equation}
\end{proposition}

\begin{proof}
We prove the slightly stronger inductive assertion in which
the full left character law is retained, but on the right
only the upper-root character laws within the first $s$
isotropic indices are required.  The stated theorem is
its special case.
The assertions for $s=0,1$ are immediate from the
empty determinant and \eqref{rt-torus-rule}.
Expand the size $s$ determinant in its last column.
The last factor in each product acts first.
By \eqref{rt-right-rule}, all terms vanish except
the last-row diagonal term and the term with row
$s-1$.  Thus its value equals that of
\[
 Q_{s-1}(u)(u-s+1+E_{ss})F-C_s(u)F.
\]
Here the abbreviated indices belong to the first
isotropic list.

The function $E_{ss}F$ retains left equivariance,
and it retains right equivariance for all upper
roots within the first $s-1$ indices, since those
generators commute with $E_{ss}$.
The induction hypothesis therefore applies to
this function.  The resulting diagonal term is
$K_{s-1}(u)(u-s+1+p_s)F_m$.
The operator $K_{s-1}(u)$ uses only
$c_1,\ldots,c_{s-1}$, so it commutes with $p_s$.

Apply Lemma~\ref{rt-cofactor-reduction} to the
other term.  The strengthened lower-root rule
in Lemma~\ref{rt-root-rules} kills every summand
$E_{sj}V_jF$ with $j<s-1$.  The remaining term
equals
\[
 \frac{c_s}{c_{s-1}}\,
       \bigl(Q_{s-2}(u)F\bigr)(b(c)m).
\]
The induction hypothesis gives exactly
\eqref{rt-K-rec}.  This proof uses only the root
laws actually needed at each stage; it does not
assume that arbitrary right derivatives preserve
all right character laws.
\end{proof}

\subsection{The exact scalar conjugation and the conjugate operator}

Define
\begin{equation}\label{rt-gauge}
 \gamma(c)=\prod_{i=1}^{r}
       |c_i|^{\,2(r-i+1)+k},
 \qquad
 \Phi(c,m)=\gamma(c)F(b(c)m).
\end{equation}
This is exactly the source's scalar radial
normalization.  Let
\[
 T_i(Y)=
 \begin{pmatrix}Y+p_i&-c_i\\ c_i^{-1}&0\end{pmatrix},
 \qquad
 \overline T_i(Y)=
 \begin{pmatrix}Y-\bar p_i&-\bar c_i\\
                         \bar c_i^{-1}&0\end{pmatrix},
\]
where entries act on functions of $c$.
Write their ordered products as
\[
 \begin{pmatrix}A_r(Y)\\ B_r(Y)\end{pmatrix}
       =T_r(Y)\cdots T_1(Y)\binom10,\qquad
 \begin{pmatrix}\overline A_r(Y)\\ \overline B_r(Y)\end{pmatrix}
       =\overline T_r(Y)\cdots\overline T_1(Y)\binom10.
\]
For $r=0$ both columns are $(1,0)^{\mathsf t}$.
Set $a(c)=\prod_i c_i$.

\begin{theorem}\label{rt-actual-lifts}
For every smooth $F$ satisfying
\eqref{rt-whittaker}, and every allowed cell point,
\begin{align}
 \gamma(c)(\mathcal A_YF)(b(c)m)
       &=a(c) A_r(Y)\Phi(c,m),\label{rt-lift-first}\\
 \gamma(c)(\overline{\mathcal A}_YF)(b(c)m)
       &=\overline{a(c)}\,\overline A_r(Y)\Phi(c,m).
                                             \label{rt-lift-second}
\end{align}
Both identities hold for every complex $Y$.
The full-group operators on the left are the
polynomial operators already constructed in
\eqref{rt-A} and \eqref{rt-Av}.
\end{theorem}

\begin{proof}
At a cell point, $h(b(c)me_j,e_a)$ is zero unless
$j=e_a$, and its nonzero value is $c_a$.
Only one scalar minor in \eqref{rt-A} survives.
Consequently
\[
 (\mathcal A_YF)(b(c)m)
      =a(c)\bigl(Q_{I,I}(u)F\bigr)(b(c)m),
             \qquad u=Y+\frac{2r+k}{2}.
\]
For $n_i=2(r-i+1)+k$,
\[
 \gamma p_i\gamma^{-1}=p_i-\frac{n_i}{2},
 \qquad u-i+1-\frac{n_i}{2}=Y.
\]
Multiplication by $c_s/c_{s-1}$ commutes with
$\gamma$.  Conjugating \eqref{rt-K-rec} therefore
gives the recurrence
\[
 A_s=(Y+p_s)A_{s-1}
                -\frac{c_s}{c_{s-1}}A_{s-2},
\]
with $A_0=1$, $A_1=Y+p_1$.
The first column of $T_s\cdots T_1$ obeys the
same recurrence, since $B_{s-1}=c_{s-1}^{-1}A_{s-2}$.
This proves \eqref{rt-lift-first}.

For the second operator,
$(b(c)m)^{-1}e_a=\bar c_a e_a^\vee$.
Thus \eqref{rt-Av} reduces to
\[
 \overline{a(c)}\,
 Q_{I^\vee,I^\vee}(w)F(b(c)m),
        \qquad w=Y-\frac{2r+k}{2}+r-1 .
\]
Simultaneously reversing both index lists of a
full Capelli determinant does not change it,
by its covariance under a permutation matrix.
Apply the proof of Proposition~\ref{rt-raw-toda}
to the reversed dual list in
Lemma~\ref{rt-root-rules}.
At the vertex corresponding to the original
index $i$, the diagonal factor is
\[
                w-(r-i)-\bar p_i,
\]
and the edge between $i$ and $i+1$ has factor
$-\bar c_{i+1}/\bar c_i$.
Conjugation by $\gamma$ changes that diagonal
factor to $Y-\bar p_i$, because
\[
 w-(r-i)+\frac{n_i}{2}=Y.
\]

For completeness, the resulting recurrence can
be read in either direction along the ordered
set of vertices $1,\ldots,r$.
Its expansion is a sum over sets of pairwise
disjoint adjacent edges.  An edge $(i,i+1)$
contributes $-\bar c_{i+1}/\bar c_i$; every
uncovered vertex $j$ contributes $Y-\bar p_j$.
All factors in a fixed term commute: an
uncovered derivative involves none of the
variables in the edge factors.  Splitting these
sets according to whether the last vertex is
uncovered or lies in its last edge gives the
recurrence for $\overline T_r\cdots\overline T_1$.
Splitting at the first vertex gives the reversed
dual recurrence.  They are therefore equal.
This proves \eqref{rt-lift-second}, including
$r=0$ and $r=1$.
\end{proof}

\paragraph{Scope.}
The two polynomial full-group operators, their
translation covariance, their action on the
coefficient class, and their exact radial
identities are now proved.  Together with
Section~\ref{rl-section}, this supplies three
concrete ingredients of the four-operator
Fourier construction.  The remaining coupled
operators involving the off-diagonal terminal
entries require their own identities.  The
passage from formal adjoints to the actual
Jacquet pairing, the uniform square-integrability
bounds for the generated class, and the final
weighted Fourier estimate are not consequences
of the two polynomial radial identities alone.


